\documentclass[12pt]{article}
\usepackage{amsmath,amssymb}

\begin{document}
\section*{{2021 AMC 12A Problems/Problem 10}}
Source: AoPS Wiki

\subsection*{Solution}
Initial Scenario Let the heights of the narrow cone and the wide cone be $h_1$ and $h_2,$ respectively. We have the following table: \[\begin{array}{cccccc} & \textbf{Base Radius} & \textbf{Height} & & \textbf{Volume} & \\ [2ex] \textbf{Narrow Cone} & 3 & h_1 & & \frac13\pi(3)^2h_1=3\pi h_1 & \\ [2ex] \textbf{Wide Cone} & 6 & h_2 & & \hspace{2mm}\frac13\pi(6)^2h_2=12\pi h_2 & \end{array}\] Equating the volumes gives $3\pi h_1=12\pi h_2,$ which simplifies to $\frac{h_1}{h_2}=4.$ Furthermore, by similar triangles: For the narrow cone, the ratio of the base radius to the height is $\frac{3}{h_1},$ which always remains constant. For the wide cone, the ratio of the base radius to the height is $\frac{6}{h_2},$ which always remains constant. Two solutions follow from here:

\end{document}
